\subsection{Архитектура предлагаемого решения}

Разработанная система представляет собой инкрементальный конвейер (pipeline) для автоматического обнаружения и типизации семантических связей между заметками в хранилище Obsidian.
Конвейер спроектирован как многоэтапный процесс: от сканирования файловой системы до записи классифицированных связей обратно в Markdown-файлы.
Ключевой особенностью архитектуры является поддержка контрольных точек (checkpoints), позволяющая возобновить работу после прерывания с того этапа, на котором она была остановлена.

Общая схема конвейера представлена на рисунке~\ref{fig:pipeline}.

\begin{figure}[H]
    \centering
    \begin{tikzpicture}[
            node distance = 6mm and 30mm,
            block/.style = {
                    rectangle, rounded corners=3pt,
                    draw=black!60, fill=blue!8,
                    text width=72mm, minimum height=9mm,
                    align=center, font=\small
                },
            ioblock/.style = {
                    rectangle, rounded corners=3pt,
                    draw=black!40, fill=gray!10,
                    text width=72mm, minimum height=8mm,
                    align=center, font=\small\itshape
                },
            sideblock/.style = {
                    rectangle, rounded corners=3pt,
                    draw=black!50, fill=orange!12,
                    text width=36mm, minimum height=8mm,
                    align=center, font=\small
                },
            arrow/.style = {-Stealth, thick, black!70},
            dasharrow/.style = {Stealth-Stealth, dashed, thick, black!50},
            label/.style = {font=\scriptsize\itshape, text=black!60}
        ]

        % Main flow
        \node[ioblock]  (vault)    {Хранилище Markdown-файлов (vault)};
        \node[block,    below=of vault]  (scan)     {1. Сканирование и обнаружение изменений\\{\scriptsize\normalfont VaultManager}};
        \node[block,    below=of scan]   (index)    {2. Разбивка на чанки и индексация\\{\scriptsize\normalfont VectorStore / LanceDB}};
        \node[block,    below=of index]  (retrieval){3. Гибридный поиск кандидатных пар\\{\scriptsize\normalfont BM25 + векторный поиск}};
        \node[block,    below=of retrieval] (rerank){4. Переранжирование кандидатов\\{\scriptsize\normalfont Cross-Encoder (reranker)}};
        \node[block,    below=of rerank] (llm)      {5. Классификация типа связи\\{\scriptsize\normalfont LLM (локально или API)}};
        \node[block,    below=of llm]    (export)   {6. Запись связей в хранилище\\{\scriptsize\normalfont ExportService}};
        \node[ioblock,  below=of export] (result)   {Обновлённые заметки с типизированными связями};

        % Side checkpoint block, aligned vertically between retrieval and llm
        \node[sideblock, right=of rerank] (meta)
        {Метаданные\\\& чекпоинты\\{\scriptsize\normalfont MetadataManager}};

        % Main arrows
        \draw[arrow] (vault)     -- node[label, right]{изменённые файлы}    (scan);
        \draw[arrow] (scan)      -- node[label, right]{новые / обновлённые файлы} (index);
        \draw[arrow] (index)     -- node[label, right]{чанки + эмбеддинги}  (retrieval);
        \draw[arrow] (retrieval) -- node[label, right]{кандидатные пары}     (rerank);
        \draw[arrow] (rerank)    -- node[label, right]{отфильтрованные пары} (llm);
        \draw[arrow] (llm)       -- node[label, right]{тип связи + reasoning}(export);
        \draw[arrow] (export)    --                                           (result);

        % Dashed arrows to/from checkpoint block
        \draw[dasharrow] (meta.west) -- (rerank.east);
        \draw[dasharrow] ([yshift=4mm]meta.west) -- ([yshift=4mm]retrieval.east);
        \draw[dasharrow] ([yshift=-4mm]meta.west) -- ([yshift=-4mm]llm.east);

    \end{tikzpicture}
    \caption{Общая схема конвейера построения графа знаний}
    \label{fig:pipeline}
\end{figure}

\textbf{Этап 1. Сканирование хранилища и обнаружение изменений (VaultManager).}
На этом этапе модуль \texttt{VaultManager} рекурсивно обходит директорию хранилища,
формируя список всех \texttt{markdown} файлов.
Для каждого файла вычисляется хеш его содержимого для иденификации изменений.
Полученный хэш сравнивается с сохраненным в базу метаданных,
которая может присутствовать в директории с предыдущих запусков.
В результате файлы разбиваются на три категории: новые, измененные и удаленные.
Наконец, новые и измененные файлы заного индексируются и записываются в векторное хранилище LanceDB,
удаленные файлы, соответственно, из него удаляются.
Данный подход позволяет экономить вычислительные ресурны, поскольку обрабатывает только изменения,
а не все хранилище целиком.


\textbf{Этап 2. Разбивка на чанки и индексация (VectorStore / LanceDB).}
Каждый новый или измененный чанк проходит процесс индексации.
Эта процедура состоит из двух частей: подготовки текстового и векторного индекса.

Текстовый индекс строится на основе лемматизированного содержимого файлов, проведенного при помощи \texttt{spacy}.
Данный подход позволяет нивелировать проблему, связанную с изменением форм слов,
в которых, например, могут быть изменены окончания или беглые гласные в корне, что особенно распространено в русском языке.

Построение векторного индекса начинается с разбиения содержимого файла на перекрывающиеся
чанки. Поддерживается две стартегии разбиения: рекурсивная посимвольная с помощью \texttt{RecursiveCharacterTextSplitter}
из \texttt{LangChain} и оконная по предложениям \texttt{SentenceWindowSplitter}, реализованныя с помощью spacy.
Далее для каждого чанка следует вычисление плотного векторного эмбеддинга при помощи модели из Sentence-Transformers.

Полученные индексы объединяются друг с другом, а также с метаданным в виде пути к файлу и хеша,
после чего сохраняются в базу данных.
LanceDB обеспечивает как векторный поиск по приближённым ближайшим соседям (ANN), так и полнотекстовый поиск (BM25).


\textbf{Этап 3. Гибридный поиск пар-кандидатов (Retrieval).}
На данном этапе для каждого чанка выполняется поиск семантически близких чанков из других файлов.
Доступно три алгоритма поиска:
\textit{broad} (векторный поиск),
\textit{strict} (поиск упоминаний сущностей в лемматизированном тексте) и
\textit{combined} (объединение обеих стратегий с дедупликацией, приоритизируются результаты \textit{strict}).
Результатом этапа является список структур \texttt{CandidatePair},
каждая из которых содержит исходный чанк, целевой чанк, пути к соответствующим файлам и оценку близости.
В случае \textit{strict} оценка всегда равна 1, в случае \textit{broad} оценка берется из LanceDB.


\textbf{Этап 4. Переранжирование кандидатов (Cross-Encoder).}
Первичный поиск зачастую возвращает большое число кандидатов, часть которых является ложноположительными.
Для снижения шума применяется двухэтапная фильтрация:
сначала выбираются \textit{top-k} результатов по оценке поиска,
затем каждая пара (чанк источника, чанк цели) подаётся на вход кросс-энкодеру (модель типа BAAI/bge-m3 / MiniLM).
Кросс-экнодер вычисляет более точную оценку семантической связности.
Пары сортируются по убыванию оценки. В финальный список попадают \textit{top-k},
у каждой их которых оценка не ниже порогового значения.

\textbf{Этап 5. Классификация типа связи (LLM).}
Прошедшие фильтрацию пары-кандидаты передаются в языковую модель для определения точного типа семантической связи.
Пары группируются по комбинации «файл-источник~--- файл-цель».
Для каждой такой группы отбираются не более \(N\) наиболее релевантных чанков для формирования контекста.
Запросы к LLM выполняются параллельными батчами, то есть в рамках одного запроса модель может классифицировать сразу несколько связей,
при этом модель обрабатывает несколько таких запросов одновременно.
Модель получает системный промпт со списком допустимых типов связей и инструкциями,
а также пользовательский промпт с пронумерованными парами файлов и их контекстом. Ответ модели возвращается в формате JSON для облегчения его обработки.
Поддерживаются как локальные модели (llama.cpp через LangChain), так и облачные провайдеры (OpenAI, Google, и другие, поддреживаемые LangChain).
После каждого батча предсказания модели сохраняются в файл, формируя чекпоинт, что позволяет возобновить классификацию в случае ошибки или краша программы. 
Данная оптимизация призвана экономить вычислительные ресурсы, убирая потребность заново делать выполненную работу.

\textbf{Этап 6. Запись связей в хранилище (ExportService).}
На заключительном этапе классифицированные связи сохраняются в одном из трех режимов:
\textit{inplace}~--- связи добавляются непосредственно в исходные Markdown-файлы под специальным заголовком;
\textit{json}~--- результаты экспортируются в JSON-файл без изменения исходных заметок;
\textit{export}~--- создаётся полная копия хранилища с вписанными связями, исходные файлы остаются нетронутыми.
Формат ссылки конфигурируется (по умолчанию~--- синтаксис Dataview: \texttt{- relation:: [[target]]}).

\textbf{Управление состоянием и отказоустойчивость (MetadataManager).}
Крайне полезным модулем, сопровождающим все этапы алгоритма, является \texttt{MetadataManager}.
Он занимается сохранением хешей файлов, конфигурации алгоритма, списока кандидатов 
и предсказаниями LLM в специальной директории \texttt{.kg\_builder/} внутри хранилища.
При следующем запуске конвейер автоматически определяет параметры, модели и этап, на котором был прерван предыдущий запуск.
Это позволяет продолжить работу с этой точки, не повторяя уже выполненных шагов.
